\begin{pastexam}[label={px:ntu-112-dsa-08}]{臺大資工所 112 學年度 第 8 題}
（複選；5 分，每個選項單獨計分，答對得 1 分、答錯倒扣 0.5 分）

\textbf{Postfix advantages:} What is/are the advantage(s) of using postfix notation for math
expressions?
\begin{enumerate}[label=(\Alph*)]
\item No need to use parentheses
\item No need to consider precedence of operators
\item Easier for human to read
\item Easier for computers to evaluate
\item More concise when compared with prefix notation
\end{enumerate}
\end{pastexam}

\begin{pastsoln}
答案：(A)、(B)、(D)。

後序式（postfix）把運算子寫在兩個運算元後面（資結書第~6.3~節）。 例子：中序
$(a + b) \times c$ 和 $a + b \times c$，後序分別是 \texttt{a b + c *} 和 \texttt{a b c * +}。
\begin{itemize}
\item \textbf{(A) 對。} 上面兩個中序式要靠括號區分，後序式不用括號就已經不同。 每個運算子作用在
它前面最近的兩個「完整的值」上，計算順序完全由寫的位置決定。
\item \textbf{(B) 對。} 計算後序式時，由左到右遇到運算子就立刻算，從來不必比較兩個運算子誰先
（資結書虛擬碼~6.7）。 優先順序只在\textbf{把中序轉成後序}時用到（資結書虛擬碼~6.8），轉好之後
就已經寫進順序裡了。
\item \textbf{(C) 錯。} 人習慣讀中序式；\texttt{a b c * + d e / -} 這種寫法人要在腦中模擬堆疊才
看得懂，不是比較好讀。
\item \textbf{(D) 對。} 用一個堆疊，由左到右掃一遍：運算元推入，運算子彈出兩個、算完推回，
$O(n)$ 時間，不需要回頭看或處理括號（資結書第~6.3.2~節）。
\item \textbf{(E) 錯。} 前序和後序都不用括號，符號（運算元和運算子）一模一樣、只是順序不同，
長度相同，沒有誰比較精簡。
\end{itemize}
\end{pastsoln}
